\documentclass[a4paper,12pt]{article}
\usepackage[T2A]{fontenc}
\usepackage[utf8]{inputenc}
\usepackage[russian]{babel}
\usepackage{amsmath,amssymb,mathtools}
\usepackage{cmbright}
\renewcommand{\familydefault}{\sfdefault}
\usepackage{icomma}
\usepackage{geometry}
\geometry{left=18mm,right=18mm,top=20mm,bottom=18mm}
\setlength{\headheight}{25pt}
\usepackage{microtype}
\usepackage{tabularx,booktabs,array,longtable}
\usepackage{enumitem}
\usepackage{hyperref}
\usepackage{parskip}
\usepackage{fancyhdr}
\usepackage{setspace}
\usepackage{xcolor}
\usepackage{needspace}

\definecolor{accent}{HTML}{176B6B}
\definecolor{darkaccent}{HTML}{0B4545}
\definecolor{textgray}{HTML}{2F3437}
\definecolor{softgray}{HTML}{6B7378}
\definecolor{linegray}{HTML}{D8DEDE}
\definecolor{softaccent}{HTML}{E8F2F2}

\hypersetup{colorlinks=true,linkcolor=darkaccent,urlcolor=darkaccent}
\setlength{\parindent}{0pt}
\onehalfspacing
\setlist[itemize]{leftmargin=6mm,itemsep=1.5mm,topsep=1mm}
\setlist[enumerate]{leftmargin=8mm,itemsep=2.2mm,topsep=1mm}
\renewcommand{\arraystretch}{1.28}

\pagestyle{fancy}
\fancyhf{}
\fancyhead[L]{\small\color{softgray}Сборник формул · тригонометрия}
\fancyhead[R]{\small\color{softgray}Ярослав Гаврилов}
\fancyfoot[C]{\small\color{softgray}\thepage}
\renewcommand{\headrulewidth}{0.3pt}
\renewcommand{\headrule}{\hbox to\headwidth{\color{linegray}\leaders\hrule height \headrulewidth\hfill}}

\newcommand{\sectionline}{\vspace{-1mm}\noindent\color{accent}\rule{\linewidth}{0.7pt}\color{textgray}\vspace{2mm}}
\newcommand{\key}[1]{\textbf{\color{darkaccent}#1}}
\newcommand{\formula}[1]{\[\displaystyle #1\]}
\newcommand{\R}{\mathbb{R}}
\newcommand{\Z}{\mathbb{Z}}

\begin{document}
\color{textgray}

\thispagestyle{empty}
\begin{center}
\vspace*{20mm}
{\Large\color{softgray}Профильная математика · ЕГЭ\par}
\vspace{8mm}
{\fontsize{29}{34}\selectfont\bfseries\color{darkaccent}Сборник формул\par}
\vspace{4mm}
{\fontsize{21}{25}\selectfont\bfseries Тригонометрия\par}
\vspace{3mm}
{\large Формулы, которые нужно узнавать и воспроизводить без подсказки\par}

\vspace{18mm}
\begin{tabular}{@{}rl@{}}
\color{softgray}Ученик: & \textbf{Ярослав Гаврилов}\\[2mm]
\color{softgray}Дата: & \textbf{5 октября 2026}\\[2mm]
\color{softgray}Режим: & \textbf{3--5 минут активного воспроизведения ежедневно}
\end{tabular}

\vfill
{\large\bfseries\color{darkaccent}Лёвин Артём Александрович\par}
\vspace{1mm}
{\normalsize эксклюзивно для Ярослава Гаврилова\par}
\vspace{14mm}
{\color{accent}\rule{0.72\linewidth}{1.2pt}}
\end{center}
\clearpage

\section*{Как работать со сборником}
\sectionline

Этот материал собран по текущей линии занятий 19.09--05.10.26. Его задача --- превратить формулы из «узнаваемых» в \key{воспроизводимые по памяти}.

\begin{enumerate}
  \item Сначала закройте правую часть страницы или формулу рукой и попробуйте восстановить запись полностью.
  \item Сверяйте не только коэффициенты, но и \key{знак}, \key{аргумент}, \key{период}, ограничения и параметр $n\in\Z$.
  \item Формулы двойного угла и основное тригонометрическое тождество учите \key{в обе стороны}.
  \item Если формула не восстановилась с первой попытки, повторите именно её в конце занятия.
  \item После формульного блока переходите к короткому применению: формула должна упрощать выражение, а не усложнять его.
\end{enumerate}

\begin{center}
\fbox{\parbox{0.88\linewidth}{\centering
\textbf{Ключевой вопрос перед преобразованием:}\\[1mm]
«Какая формула прямо сейчас использует данные задачи и делает выражение проще?»
}}
\end{center}

\section*{1. Единичная окружность и табличные значения}
\sectionline

\subsection*{Координатный смысл}

Если точка $M$ на единичной окружности соответствует углу $\alpha$, то
\[
\boxed{M(\cos\alpha;\,\sin\alpha)}.
\]
Отсюда
\[
\boxed{x_M=\cos\alpha},\qquad \boxed{y_M=\sin\alpha}.
\]

Полный оборот:
\[
\boxed{2\pi=360^\circ}.
\]
Углы $\alpha$ и $\alpha+2\pi n$, $n\in\Z$, задают одну и ту же точку окружности.

\subsection*{Табличные углы}

\begin{center}
\begin{tabularx}{0.98\linewidth}{>{\centering\arraybackslash}X*{5}{>{\centering\arraybackslash}X}}
\toprule
\textbf{Угол} & $0^\circ$ & $30^\circ$ & $45^\circ$ & $60^\circ$ & $90^\circ$\\
\midrule
\textbf{Радианы} & $0$ & $\dfrac\pi6$ & $\dfrac\pi4$ & $\dfrac\pi3$ & $\dfrac\pi2$\\[2mm]
$\sin\alpha$ & $0$ & $\dfrac12$ & $\dfrac{\sqrt2}{2}$ & $\dfrac{\sqrt3}{2}$ & $1$\\[2mm]
$\cos\alpha$ & $1$ & $\dfrac{\sqrt3}{2}$ & $\dfrac{\sqrt2}{2}$ & $\dfrac12$ & $0$\\
\bottomrule
\end{tabularx}
\end{center}

Быстрое восстановление таблицы:
\[
\sin:\quad \frac{\sqrt0}{2},\ \frac{\sqrt1}{2},\ \frac{\sqrt2}{2},\ \frac{\sqrt3}{2},\ \frac{\sqrt4}{2},
\]
\[
\cos:\quad \frac{\sqrt4}{2},\ \frac{\sqrt3}{2},\ \frac{\sqrt2}{2},\ \frac{\sqrt1}{2},\ \frac{\sqrt0}{2}.
\]

\subsection*{Знаки по четвертям}

\begin{center}
\begin{tabular}{c|c|c|c}
\toprule
\textbf{Четверть} & $\sin$ & $\cos$ & $\tg$\\
\midrule
I   & $+$ & $+$ & $+$\\
II  & $+$ & $-$ & $-$\\
III & $-$ & $-$ & $+$\\
IV  & $-$ & $+$ & $-$\\
\bottomrule
\end{tabular}
\end{center}

\key{Правило знака:} после извлечения корня сначала пишите $\pm$. Например,
\[
\cos^2\alpha=c\quad\Longrightarrow\quad \cos\alpha=\pm\sqrt c.
\]
Нужный знак выбирается по интервалу или четверти.

\section*{2. Основные тождества}
\sectionline

\subsection*{Главное тождество}
\[
\boxed{\sin^2\alpha+\cos^2\alpha=1}.
\]

Его нужно узнавать в обе стороны:
\[
\sin^2\alpha+\cos^2\alpha\longrightarrow1,
\qquad
1\longrightarrow\sin^2\alpha+\cos^2\alpha.
\]

\subsection*{Тангенс и котангенс}
\[
\boxed{\tg\alpha=\frac{\sin\alpha}{\cos\alpha}},\qquad \cos\alpha\ne0,
\]
\[
\boxed{\ctg\alpha=\frac{\cos\alpha}{\sin\alpha}=\frac1{\tg\alpha}},
\qquad \sin\alpha\ne0.
\]

Если известен тангенс, полезно сразу связать синус и косинус. Например,
\[
\tg\alpha=-2{,}5
\quad\Longrightarrow\quad
\sin\alpha=-2{,}5\cos\alpha.
\]

\section*{3. Формулы суммы и разности}
\sectionline

\subsection*{Синус}
\[
\boxed{\sin(\alpha+\beta)=\sin\alpha\cos\beta+\cos\alpha\sin\beta},
\]
\[
\boxed{\sin(\alpha-\beta)=\sin\alpha\cos\beta-\cos\alpha\sin\beta}.
\]

\subsection*{Косинус}
\[
\boxed{\cos(\alpha+\beta)=\cos\alpha\cos\beta-\sin\alpha\sin\beta},
\]
\[
\boxed{\cos(\alpha-\beta)=\cos\alpha\cos\beta+\sin\alpha\sin\beta}.
\]

\key{Контроль знаков:} у синуса знак внутри скобок сохраняется; у косинуса знак между произведениями противоположен знаку внутри скобок.

Два опорных следствия, которые встречались в занятиях:
\[
\boxed{\sin\left(\frac\pi2+x\right)=\cos x},
\qquad
\boxed{\cos\left(\frac\pi2-x\right)=\sin x}.
\]

\section*{4. Формулы двойного угла}
\sectionline

\subsection*{Синус двойного угла}
\[
\boxed{\sin2\alpha=2\sin\alpha\cos\alpha}.
\]

Обратная форма:
\[
\boxed{\sin\alpha\cos\alpha=\frac12\sin2\alpha}.
\]

Если аргументы имеют вид $u$ и $u/2$:
\[
\boxed{\sin u=2\sin\frac u2\cos\frac u2}.
\]

\subsection*{Косинус двойного угла}
\[
\boxed{\cos2\alpha=\cos^2\alpha-\sin^2\alpha},
\]
\[
\boxed{\cos2\alpha=2\cos^2\alpha-1},
\]
\[
\boxed{\cos2\alpha=1-2\sin^2\alpha}.
\]

\key{Как выбирать форму:}
\begin{itemize}
  \item дан $\sin\alpha$ $\Rightarrow$ чаще всего используйте $\cos2\alpha=1-2\sin^2\alpha$;
  \item дан $\cos\alpha$ $\Rightarrow$ чаще всего используйте $\cos2\alpha=2\cos^2\alpha-1$;
  \item присутствуют и $\sin\alpha$, и $\cos\alpha$ $\Rightarrow$ проверьте форму $\cos^2\alpha-\sin^2\alpha$.
\end{itemize}

\begin{center}
\fbox{\parbox{0.88\linewidth}{\centering
Если рядом встречаются аргументы $x$ и $2x$, $3x$ и $6x$, сначала проверяйте формулы двойного угла.
}}
\end{center}

\section*{5. Формулы приведения: правило вместо перегруженной таблицы}
\sectionline

Для выражений вида
\[
f\!\left(k\frac\pi2\pm\alpha\right)
\]
работайте в два шага.

\begin{enumerate}
  \item \key{Меняется ли функция?} При нечётном числе четвертей
  $\left(\frac\pi2,\frac{3\pi}{2},\frac{5\pi}{2},\ldots\right)$ функция меняется на кофункцию:
  \[
  \sin\leftrightarrow\cos,
  \qquad
  \tg\leftrightarrow\ctg.
  \]
  \item \key{Какой знак?} Определяйте знак по \emph{исходной} функции в той четверти, куда попадает угол.
\end{enumerate}

Ключевые преобразования из текущего блока:
\[
\boxed{\sin\left(\frac{7\pi}{2}-\alpha\right)=-\cos\alpha},
\]
\[
\boxed{\tg\left(\alpha+\frac{5\pi}{2}\right)=-\ctg\alpha}.
\]

\key{Типичная ошибка:} определять знак по функции, которая получилась после приведения. Знак проверяется по исходной функции.

\section*{6. Периоды, области определения и значения}
\sectionline

\subsection*{Базовые периоды}
\[
\boxed{T_{\sin x}=T_{\cos x}=2\pi},
\qquad
\boxed{T_{\tg x}=T_{\ctg x}=\pi}.
\]

Для аргумента $kx$:
\[
\boxed{T_{\sin(kx)}=T_{\cos(kx)}=\frac{2\pi}{|k|}},
\]
\[
\boxed{T_{\tg(kx)}=T_{\ctg(kx)}=\frac{\pi}{|k|}}.
\]

\subsection*{Область определения и значения}
\[
\boxed{D(\sin x)=D(\cos x)=\R},
\qquad
\boxed{-1\le\sin x\le1},
\qquad
\boxed{-1\le\cos x\le1}.
\]

\[
\boxed{D(\tg x)=\R\setminus\left\{\frac\pi2+\pi n\right\}},
\qquad n\in\Z,
\]
\[
\boxed{D(\ctg x)=\R\setminus\{\pi n\}},
\qquad n\in\Z.
\]

\[
\boxed{E(\tg x)=E(\ctg x)=\R}.
\]

Для функций
\[
y=a\sin x+b
\qquad\text{или}\qquad
y=a\cos x+b
\]
область значений:
\[
\boxed{b-|a|\le y\le b+|a|}.
\]

\section*{7. Арксинус и арккосинус}
\sectionline

Для $a\in[-1;1]$:
\[
\boxed{\arcsin a\in\left[-\frac\pi2;\frac\pi2\right]},
\qquad
\boxed{\sin(\arcsin a)=a},
\]
\[
\boxed{\arccos a\in[0;\pi]},
\qquad
\boxed{\cos(\arccos a)=a}.
\]

Опорные значения:
\[
\arcsin\frac12=\frac\pi6,
\qquad
\arcsin\left(-\frac12\right)=-\frac\pi6,
\]
\[
\arccos\frac12=\frac\pi3,
\qquad
\arccos\left(-\frac12\right)=\frac{2\pi}{3}.
\]

\key{Важно:} $\arcsin a$ и $\arccos a$ --- главные значения. Полное множество решений требует периодических серий.

\section*{8. Простейшие тригонометрические уравнения}
\sectionline

Во всех формулах ниже $n\in\Z$.

\subsection*{Синус}

Если $|a|\le1$:
\[
\boxed{x=(-1)^n\arcsin a+\pi n}.
\]

Эквивалентная запись двумя сериями:
\[
\boxed{
\begin{aligned}
x&=\arcsin a+2\pi n,\\
x&=\pi-\arcsin a+2\pi n.
\end{aligned}}
\]

Частные случаи:
\[
\boxed{\sin x=0\Rightarrow x=\pi n},
\]
\[
\boxed{\sin x=1\Rightarrow x=\frac\pi2+2\pi n},
\]
\[
\boxed{\sin x=-1\Rightarrow x=-\frac\pi2+2\pi n}.
\]

\subsection*{Косинус}

Если $|a|\le1$:
\[
\boxed{x=\pm\arccos a+2\pi n}.
\]

Частные случаи:
\[
\boxed{\cos x=0\Rightarrow x=\frac\pi2+\pi n},
\]
\[
\boxed{\cos x=1\Rightarrow x=2\pi n},
\]
\[
\boxed{\cos x=-1\Rightarrow x=\pi+2\pi n}.
\]

\subsection*{Тангенс}
\[
\boxed{\tg x=a\Rightarrow x=\arctg a+\pi n}.
\]

Например,
\[
\tg x=-1\Rightarrow x=-\frac\pi4+\pi n.
\]

\subsection*{Если аргумент сложнее, чем $x$}

Сначала решайте уравнение относительно \key{всего аргумента}, затем выражайте $x$. Например:
\[
\sin3x=1
\Rightarrow
3x=\frac\pi2+2\pi n
\Rightarrow
\boxed{x=\frac\pi6+\frac{2\pi n}{3}}.
\]

\section*{9. Отбор корней на промежутке}
\sectionline
\enlargethispage{3\baselineskip}

Если серия имеет вид
\[
x=x_0+2\pi n,
\qquad n\in\Z,
\]
а требуется отбор на $[a;b]$, подставьте серию в двойное неравенство:
\[
\boxed{a\le x_0+2\pi n\le b}.
\]

Далее найдите целые значения $n$ и вернитесь к $x$.

Для серий с тангенсом шаг равен $\pi$, поэтому подставляется соответствующая формула
\[
x=x_0+\pi n.
\]

\key{Контроль:} период в серии корней берётся из конкретной функции. Для $\tg$ используется $\pi$, для $\sin$ и $\cos$ обычно $2\pi$; в частных случаях период серии может сокращаться.

\clearpage
\section*{10. Минимум для мгновенного воспроизведения}
\sectionline

\begingroup
\small
\setstretch{1.0}
\renewcommand{\arraystretch}{1.12}
\begin{center}
\begin{tabularx}{0.98\linewidth}{>{\bfseries\color{darkaccent}}p{40mm}X}
\toprule
Координаты окружности & $M(\cos\alpha;\sin\alpha)$\\[1mm]
Основное тождество & $\sin^2\alpha+\cos^2\alpha=1$\\[1mm]
Тангенс & $\displaystyle\tg\alpha=\frac{\sin\alpha}{\cos\alpha}$\\[2mm]
Котангенс & $\displaystyle\ctg\alpha=\frac{\cos\alpha}{\sin\alpha}=\frac1{\tg\alpha}$\\[2mm]
Синус суммы/разности & $\sin(\alpha\pm\beta)=\sin\alpha\cos\beta\pm\cos\alpha\sin\beta$\\[2mm]
Косинус суммы/разности & $\cos(\alpha\pm\beta)=\cos\alpha\cos\beta\mp\sin\alpha\sin\beta$\\[2mm]
Синус двойного угла & $\sin2\alpha=2\sin\alpha\cos\alpha$\\[1mm]
Косинус двойного угла & $\cos2\alpha=\cos^2\alpha-\sin^2\alpha=2\cos^2\alpha-1=1-2\sin^2\alpha$\\[2mm]
Периоды & $T_{\sin,\cos}=2\pi$, $T_{\tg,\ctg}=\pi$\\[1mm]
Период с $kx$ & $\displaystyle T_{\sin(kx),\cos(kx)}=\frac{2\pi}{|k|}$, $\displaystyle T_{\tg(kx),\ctg(kx)}=\frac{\pi}{|k|}$\\[2mm]
Область значений $a\sin x+b$ & $b-|a|\le y\le b+|a|$\\[1mm]
$\sin x=a$ & $x=(-1)^n\arcsin a+\pi n$\\[1mm]
$\cos x=a$ & $x=\pm\arccos a+2\pi n$\\[1mm]
$\tg x=a$ & $x=\arctg a+\pi n$\\
\bottomrule
\end{tabularx}
\end{center}

\vspace{2mm}
\begin{center}
\fbox{\parbox{0.9\linewidth}{
\textbf{Формулы приведения:} для $f\!\left(k\frac\pi2\pm\alpha\right)$ сначала определите, меняется ли функция на кофункцию; затем поставьте знак по \emph{исходной} функции и четверти.
}}
\end{center}
\endgroup

\clearpage
\section*{Тренировочный блок --- Путь А}
\sectionline

\textbf{10 заданий на активное воспроизведение.} Выполняйте без подсматривания. Здесь проверяется именно память формул, поэтому сначала восстановите запись полностью, затем откройте ключ.

\begin{enumerate}
  \item Восстановите основное тригонометрическое тождество:
  \[
  \sin^2\alpha+\underline{\hspace{28mm}}=\underline{\hspace{16mm}}.
  \]

  \item Восстановите определения тангенса и котангенса:
  \[
  \tg\alpha=\frac{\underline{\hspace{18mm}}}{\underline{\hspace{18mm}}},
  \qquad
  \ctg\alpha=\frac{\underline{\hspace{18mm}}}{\underline{\hspace{18mm}}}.
  \]

  \item Восстановите две формулы синуса суммы и разности:
  \[
  \sin(\alpha+\beta)=\underline{\hspace{95mm}},
  \]
  \[
  \sin(\alpha-\beta)=\underline{\hspace{95mm}}.
  \]

  \item Восстановите компактную запись формул косинуса суммы и разности:
  \[
  \cos(\alpha\pm\beta)=\underline{\hspace{105mm}}.
  \]

  \item Восстановите формулу синуса двойного угла и её обратную форму:
  \[
  \sin2\alpha=\underline{\hspace{55mm}},
  \qquad
  \sin\alpha\cos\alpha=\underline{\hspace{55mm}}.
  \]

  \item Запишите \textbf{три} равносильные формы для $\cos2\alpha$.
  \vspace{18mm}

  \item Запишите наименьшие положительные периоды функций $\sin(kx)$, $\cos(kx)$, $\tg(kx)$ и $\ctg(kx)$.
  \vspace{18mm}

  \item Восстановите два ключевых примера формул приведения из текущего блока:
  \[
  \sin\left(\frac{7\pi}{2}-\alpha\right)=\underline{\hspace{35mm}},
  \qquad
  \tg\left(\alpha+\frac{5\pi}{2}\right)=\underline{\hspace{35mm}}.
  \]

  \item Запишите общие решения:
  \[
  \sin x=a:\quad \underline{\hspace{95mm}},
  \]
  \[
  \cos x=a:\quad \underline{\hspace{95mm}},
  \]
  \[
  \tg x=a:\quad \underline{\hspace{95mm}}.
  \]

  \item Восстановите формулу области значений для $y=a\sin x+b$ и $y=a\cos x+b$:
  \[
  \underline{\hspace{105mm}}.
  \]
\end{enumerate}

\clearpage
\section*{Ключ к тренировочному блоку}
\sectionline

\begin{enumerate}
  \item $\sin^2\alpha+\cos^2\alpha=1$.

  \item $\displaystyle\tg\alpha=\frac{\sin\alpha}{\cos\alpha}$, \quad
  $\displaystyle\ctg\alpha=\frac{\cos\alpha}{\sin\alpha}$.

  \item
  $\sin(\alpha+\beta)=\sin\alpha\cos\beta+\cos\alpha\sin\beta$;\\
  $\sin(\alpha-\beta)=\sin\alpha\cos\beta-\cos\alpha\sin\beta$.

  \item
  $\cos(\alpha\pm\beta)=\cos\alpha\cos\beta\mp\sin\alpha\sin\beta$.

  \item
  $\sin2\alpha=2\sin\alpha\cos\alpha$; \quad
  $\displaystyle\sin\alpha\cos\alpha=\frac12\sin2\alpha$.

  \item
  $\cos2\alpha=\cos^2\alpha-\sin^2\alpha=2\cos^2\alpha-1=1-2\sin^2\alpha$.

  \item
  $\displaystyle T_{\sin(kx)}=T_{\cos(kx)}=\frac{2\pi}{|k|}$; \quad
  $\displaystyle T_{\tg(kx)}=T_{\ctg(kx)}=\frac{\pi}{|k|}$.

  \item
  $\displaystyle\sin\left(\frac{7\pi}{2}-\alpha\right)=-\cos\alpha$; \quad
  $\displaystyle\tg\left(\alpha+\frac{5\pi}{2}\right)=-\ctg\alpha$.

  \item
  $\displaystyle\sin x=a\Rightarrow x=(-1)^n\arcsin a+\pi n$, $|a|\le1$;\\
  $\displaystyle\cos x=a\Rightarrow x=\pm\arccos a+2\pi n$, $|a|\le1$;\\
  $\displaystyle\tg x=a\Rightarrow x=\arctg a+\pi n$, \quad $n\in\Z$.

  \item $\boxed{b-|a|\le y\le b+|a|}$.
\end{enumerate}

\vfill
\begin{center}
{\color{accent}\rule{0.56\linewidth}{0.9pt}}\\[3mm]
{\small\color{softgray}Лёвин Артём Александрович · эксклюзивно для Ярослава Гаврилова}
\end{center}

\end{document}
